\documentclass{article}

\newif\ifdraft
\usepackage{amsmath}
\usepackage{graphicx}
\usepackage{hyperref}
\usepackage{mypackage}

% Standard single-argument macros.
\newcommand{\highlight}[1]{\textbf{#1}}
\newcommand{\norm}[1]{\left\lVert#1\right\rVert}

\csname ifdraft\endcsname\else
% Directly recursive; found by the `recursion` query.
\newcommand{\repeat}[1]{#1\repeat{#1}}
\fi
\begin{document}

\def\DocName{Document}
\title{A Sample \DocName{} for sa\TeX}
\author{satex Test Suite}
\maketitle

\section{Introduction}
\label{sec:intro}
\ifdraft
This document exercises \highlight{satex}, a static analyser for \LaTeX{}.
See \autoref{sec:method} for the methodology and \autoref{sec:results} for
the results.
\fi

\section{Methodology}
\label{sec:method}

We define the loss function as follows.

\begin{equation}
  \label{eq:loss}
  \mathcal{L} = \sum_{i=1}^{N} \norm{x_i - \hat{x}_i}^2
\end{equation}

\begin{figure}[ht]
  \centering
  \rule{0.4\linewidth}{0.3pt}
  \caption{Placeholder figure for the satex sample.}
  \label{fig:main}
\end{figure}

\section{Results}
\label{sec:results}

\autoref{eq:loss} defines the objective; \autoref{fig:main} illustrates the
model architecture.  A summary of results appears in \autoref{tab:results}.

\begin{table}[ht]
  \centering
  \caption{Summary of results.}
  \label{tab:results}
  \begin{tabular}{lc}
    Method & Score \\ \hline
    Ours   & 0.92  \\
  \end{tabular}
\end{table}

\section{Conclusion}
\label{sec:conclusion}

\highlight{satex} correctly analyses this document.
See \autoref{sec:intro} for motivation and \autoref{sec:results} for
the quantitative evaluation.

\end{document}
